
\begin{figure}[htbp]
  \centering
  \includegraphics[width=0.86\textwidth]{pic/fig1-1.png}
  \caption{研究内容与技术路线图}
\end{figure}











\section{机场牵引车辆系统建模与问题理论基础}
\subsection{机场牵引作业场景抽象}
机场牵引作业通常发生在停机位、廊桥附近和机坪连接滑行通道区域。与
普通道路车辆行驶相比，牵引车的运行速度很低，但控制难度并不低：车辆要
在有限空间内靠近飞机前起落架，既要避免车体、机翼和地面保障设备之间发
生干涉，又要保证车头方向、横向位置和停车距离满足后续机械连接要求。对
于本课题中的双桥转向牵引车而言，前后转向桥能够提高机动性，但也使车辆
运动模式和控制输入之间的关系更加复杂。
结合本文任务，可把自动牵引对接看成几个连续但控制重点不同的阶段：
任务初始化时确认目标飞机；初始接近阶段从出发位姿驶向预对接区域；姿态
对准阶段重点压低横向误差和航向误差；微动对接阶段以较小速度修正剩余偏
差；停车保持阶段在距离、姿态和安全阈值满足后发布零轮速和转向回中命
令。这样划分后，后续状态机和控制参数的设置更清楚。
由此可见，飞机推出不是一条轨迹从起点跟到终点这么简单，而是目标识
别、车辆运动学约束、安全边界和模式切换共同作用的状态演化过程[17-21]。本
章只先处理牵引车本体在机坪平面内如何运动的问题，即建立位姿状态、曲率
关系和输入约束；飞机辅助识别、局部安全判断和最终对接控制将在第三章统
一展开。
\subsection{坐标系与状态变量定义}
为统一描述牵引车在机坪平面内的运动状态，建立全局惯性坐标系 𝑂 −
𝑋𝑌与车体坐标系 𝐶 − 𝑥𝑦。其中，全局坐标系固定于地面，主要用于描述车辆
在场景中的绝对位置；车体坐标系原点取在车辆参考点 𝐶处，该点可取车辆几
何中心、质心或控制参考中心，𝑥 轴沿车辆纵向前进方向，𝑦 轴沿车辆横向方
向。
定义车辆位姿状态为𝐪 = [𝑥    𝑦   𝜃]T   。其中，𝑥 与 𝑦分别表示车辆参考
点在全局坐标系中的位置坐标，𝜃 表示车体纵向轴相对于全局 𝑋轴的夹角，即
车辆航向角。进一步地，为描述车辆运动过程中的速度状态，引入线速度 𝑣作
为系统运动输入或状态扩展量，则可将系统状态写为𝐱 = [𝑥           𝑦   𝜃   𝑣]T 。式
中𝑣 表示车辆沿车体纵向的速度分量。该状态定义方式具有两点作用：一方
面，(𝑥, 𝑦, 𝜃) 可直接用于描述车辆轨迹；另一方面，𝑣 可与转向输入共同决定
系统状态的实时演化。







对于本课题所研究的双桥转向牵引车，前桥和后桥均具有独立转向能力，
因此进一步定义前桥转角与后桥转角分别为 𝛿𝑓 与 𝛿𝑟 。若将系统控制输入写为
𝐮 = [𝑣   𝛿𝑓   𝛿𝑟 ]T ，则后续运动学模型即可统一写成关于 𝐱与 𝐮的非线性状态
方程。
\subsection{车辆运动学模型建立}
\subsubsection{非完整约束条件}
牵引车在机坪低速运行过程中，车轮主要满足纯滚动约束，即轮胎可沿轮
平面方向滚动，但不能在垂直于轮平面的方向上发生理想横向滑移。因此，车
辆系统属于典型非完整约束系统。对于车体参考点 𝐶，其速度在车体横向方向
上的投影应为零，可写为𝑦̇ cos 𝜃 − 𝑥̇ sin 𝜃 = 0。该式的物理意义在于：车辆不
能瞬时“侧向平移”，其平面运动只能由纵向速度与航向变化共同形成。这一
约束是移动机器人与地面车辆运动学建模中的基本条件，也是后续推导轨迹方
程的前提[17,19]。
\subsubsection{单轨等效模型}
由于前桥左右轮、后桥左右轮在同一转向控制下可近似视为对称运动，因
此可采用单轨等效方法，将实际四轮结构抽象为前后两虚拟轮[19]。设前、后桥
中心点分别位于车体参考点前后，车辆轴距为 𝐿。在低速、小侧偏工况下，车
辆瞬时旋转中心满足几何关系，车辆平面运动可表示为（1）
𝑥̇ = 𝑣cos 𝜃
\{ 𝑦̇ = 𝑣sin 𝜃                     （1）
𝜃̇ = 𝜔

式中，𝜔 为车辆偏航角速度。前两式给出了车辆质点速度在全局坐标系中
的分解形式，第三式给出了车辆姿态随时间的演化规律。车辆轨迹由两个因素
共同决定：一是沿当前朝向的平移速度 𝑣，二是决定航向变化速率的角速度
𝜔。
\subsubsection{双桥转向角速度表达式}
对于普通前轮转向车辆，角速度通常由单一前轮转角决定；而对于双桥转
向车辆，前桥和后桥转向均会对瞬时旋转中心位置产生影响。根据单轨几何关
系，前桥与后桥分别对应的曲率贡献为 tan 𝛿𝑓 /𝐿与 tan 𝛿𝑟 /𝐿。在统一参考方向
𝑣
下，车辆偏航角速度可表示为：𝜃̇ = 𝜔 = 𝐿 (tan 𝛿𝑓 − tan 𝛿𝑟 ) 。𝛿𝑓 为前桥转
角，𝛿𝑟 为后桥转角，𝐿 为前后桥中心距。该式表明，前桥转向与后桥转向对车
辆角速度的作用方向相反。当后桥不转向，即 𝛿𝑟 = 0时，模型退化为普通前轮






转向模型；当前后桥反向转向时，tan 𝛿𝑓 − tan 𝛿𝑟 的绝对值增大，车辆角速度
提高，转弯半径减小；当前后桥同向转向时，两者效应相互抵消，车辆偏航角
速度下降，车辆运动趋于“蟹行”模式。
将上述关系代入状态方程，可得双桥转向牵引车的运动学模型（2）
𝑥̇ = 𝑣cos 𝜃
\{       𝑦̇ = 𝑣sin 𝜃            （2）
𝑣
𝜃̇ = 𝐿 (tan 𝛿𝑓 − tan 𝛿𝑟 )



该模型是本课题后续路径规划、轨迹跟踪与对接控制设计的基础模型
[17,19,28]
。
\subsection{转弯半径与曲率模型}
车辆轨迹几何特性通常由曲率或转弯半径刻画。设车辆瞬时转弯半径为
𝑅，则有
𝑣
𝜔=𝑅                （3）


结合上一节所得角速度表达式，可得双桥转向牵引车的等效转弯半径为
𝐿
𝑅=                      （4）
tan 𝛿𝑓 −tan 𝛿𝑟


若进一步定义轨迹曲率 𝜅为转弯半径的倒数，则有
1     tan 𝛿𝑓 −tan 𝛿𝑟
𝜅=𝑅=                         （5）
𝐿


由式（4）和式（5）可以看出，前后桥转角最终会反映到轨迹弯曲程度
上。转弯半径描述车辆完成转向所需的空间，曲率则描述航向变化的快慢。对
本文的机坪对接场景而言，车辆能否在较短距离内调整车头方向，实际取决于
当前转角组合所能提供的曲率范围，以及控制器是否能让该曲率平稳变化
[19,28]。𝑅𝑅𝜅𝜅
还需要注意，公式中的极限情况只是几何推导结果。在真实机构或 Gazebo
模型中，前后桥转角会受到关节限位、控制器响应和模型稳定性的共同限制。
曲率不能简单地无限增大，否则会带来轮组抖动、车身摆动或轨迹不连续。因
此，后续规划和控制只能使用车辆实际能够执行的曲率范围。tan 𝛿𝑓 −
tan 𝛿𝑟 = 0𝑅 → ∞ ∣ tan 𝛿𝑓 − tan 𝛿𝑟 ∣
\subsection{典型转向模式分析}






双桥转向车辆较普通单桥转向车辆具有更丰富的运动模式。为明确不同转
向策略对车辆机动性的影响，分别讨论以下三种典型工况。


\subsubsection{前桥单独转向模式}
当后桥不参与转向，即𝛿𝑟 = 0时，车辆运动学模型退化为
𝑣
𝜃̇ = 𝐿 tan 𝛿𝑓                （6）

此时车辆行为与传统前轮转向车辆基本一致（图 2-1）。该模式实现简单、
控制逻辑清晰，但在对接空间受限的场景下，其最小转弯半径较大，不利于快
速完成姿态调整。




\begin{figure}[htbp]
  \centering
  \includegraphics[width=0.80\textwidth]{pic/fig2-1.png}
  \caption{传统转向模式}
\end{figure}


\subsubsection{前后桥反向转向模式}
当后桥转角与前桥转角方向相反，且幅值近似相等，即𝛿𝑟 = −𝛿𝑓 代入角速
表达式可得
𝑣                         2𝑣
𝜃̇ = 𝐿 (tan 𝛿𝑓 − tan (−𝛿𝑓 )) = 𝐿 tan 𝛿𝑓   （7）

此时车辆偏航角速度约为单桥转向工况的两倍，转弯半径显著减小（图 2-
2）
，适用于牵引车在狭小机坪空间内快速调整姿态、靠近飞机前轮或完成局部
避障。该模式是本课题中实现高机动性接近控制的重要运动模式。











\begin{figure}[htbp]
  \centering
  \includegraphics[width=0.80\textwidth]{pic/fig2-2.png}
  \caption{前后桥反向转向模式}
\end{figure}


\subsubsection{前后桥同向转向模式}
𝑣
当𝛿𝑟 = 𝛿𝑓 时，有𝜃̇ = 𝐿 (tan 𝛿𝑓 − tan 𝛿𝑓 ) = 0
当车辆偏航角速度为零，此时车辆整体保持原有航向并沿车轮切线方向斜
向平移，表现为 “蟹行”运动（图 2-3）。该模式并不适合大幅姿态修正，但
对于微小横向误差修正、精细对位和低速贴近控制具有较高实用价值。
通过上述分析可知，双桥转向结构使牵引车不再局限于单一的“直行—转
弯”运动，而具备在姿态快速修正与横向精细调节之间切换的能力，这也是其
适用于飞机自动对接场景的重要原因。




\begin{figure}[htbp]
  \centering
  \includegraphics[width=0.80\textwidth]{pic/fig2-3.png}
  \caption{前后桥同向转向模式}
\end{figure}


\subsection{速度、转角与曲率约束建模}
运动学模型给出了系统状态演化规律，但在实际推引作业中，车辆运动必
须满足结构约束与运行约束。因此，需要进一步建立输入与状态约束。
\subsubsection{速度约束}
牵引车在机场作业环境中通常处于低速运行状态，故其线速度应满足0 ≤
𝑣 ≤ 𝑣max 。






其中，𝑣max 为车辆允许的最大作业速度。该约束保证车辆在接近、对接和
牵引过程中具有足够的可控性，同时减小碰撞风险。
\subsubsection{转向角约束}
由于前后桥转向机构存在机械极限，故应满足
∣ 𝛿𝑓 ∣≤ 𝛿𝑓,max
\{                                                             （8）
∣ 𝛿𝑟 ∣≤ 𝛿𝑟,max

其中，𝛿𝑓,max 与 𝛿𝑟,max 分别为前桥与后桥最大允许转角。该约束直接限制
车辆最大可实现曲率，也是最小转弯半径存在下界的原因。
\subsubsection{曲率约束}
由转向角约束可进一步得到曲率约束∣ 𝜅 ∣≤ 𝜅max 。其中，
tan 𝛿𝑓,max −tan (−𝛿𝑟,max )
𝜅max =                                               （9）
𝐿


当车辆采用对称反向转向策略时，上式取到较大值。曲率约束在路径规划
中具有重要作用，因为若规划轨迹的局部曲率超过车辆可实现上限，则该轨迹
在物理上不可执行[20-21,28]。
\subsection{车辆动力学简化与驱动关系}
本课题以低速自主牵引控制为研究重点，后续控制层主要关注位姿演化与
轨迹误差收敛，因此在第二章中不引入完整轮胎力学模型，而仅保留纵向驱动
的简化动力学关系[19]。设车辆等效质量为 𝑚，驱动力为 𝐹，纵向加速度为 𝑎，
则有𝐹 = 𝑚𝑎 ；若用速度导数表示，则有𝑎 = 𝑣̇ 。
从而可得
𝐹 = 𝑚𝑣̇                                  （10）

式（10）用于说明驱动输入与速度变化之间的关系。本文没有把完整轮胎
力学放入模型，而是通过车轮速度控制器实现纵向驱动：控制层给出期望线速
度，仿真层再把它转换为车轮角速度。对于本文低速对接验证而言，这种简化
已经能够支撑车辆前进、减速和停车；若后续研究载荷变化或制动响应，再补
充更完整的动力学模型。𝐹
\subsection{状态方程统一表达}
综合车辆位姿状态、双桥转向输入与速度演化关系，可将系统模型统一表
示为非线性状态方程𝐱̇ = 𝑓(𝐱, 𝐮)。
其中，𝐱 = [𝑥   𝑦   𝜃   𝑣]T ，𝐮 = [𝑣     𝛿𝑓   𝛿𝑟 ]T 或𝐮 = [𝑎      𝛿𝑓   𝛿𝑟 ]T ；








若采用速度直接控制形式，则状态方程可写为
𝑥̇ = 𝑣cos 𝜃
\{       𝑦̇ = 𝑣sin 𝜃              （11）
𝑣
𝜃̇ = 𝐿 (tan 𝛿𝑓 − tan 𝛿𝑟 )


若进一步考虑纵向加速度，则可扩展为
𝑥̇ = 𝑣cos 𝜃
𝑦̇ = 𝑣sin 𝜃
𝑣                                  （12）
𝜃̇ = 𝐿 (tan 𝛿𝑓 − tan 𝛿𝑟 )
\{          𝑣̇ = 𝑎

式（11）适合直接速度控制和基本路径跟踪，变量较少，便于在 ROS 节点
中实时计算；式（12）多考虑了纵向加速度，后续可用于速度平滑、模型预测
控制或轨迹优化。本文当前重点是验证牵引车能否完成接近、对准和停车，因
此主要采用较简洁的速度控制形式。
\subsection{仿真实现中的参数映射关系}
在 ROS-Gazebo 联合仿真中，公式模型并不会以连续方程直接运行。车辆动
作需要经过 URDF 关节、controller 配置和 ROS 话题发布才能体现到模型上
[11-16]。因此，本节需要把理论中的速度、转角等变量对应到 Gazebo 中的关
节位置命令和轮速命令，保证论文模型与仿真执行对象保持一致。
具体到本文模型，前桥转角和后桥转角分别发送给两个转向关节的位置控
制器；车辆线速度则换算为四个驱动轮的角速度命令。设车轮半径为，则车辆
线速度与轮速之间满足𝛿𝑓 𝛿𝑟 𝑣𝑟𝑤 𝜔𝑤
𝑣
𝜔𝑤 = 𝑟                   （13）
𝑤



若四个驱动轮采用相同滚动半径并同步驱动，则其控制目标可统一写为
𝑣
𝜔𝑓𝑙 = 𝜔𝑓𝑟 = 𝜔𝑟𝑙 = 𝜔𝑟𝑟 = 𝑟         （14）
𝑤



其中，四个角速度分别对应前左、前右、后左、后右车轮。该映射把第二
章中的线速度变量接到了 Gazebo 中的 velocity\_controller 话题上，是后续通
过“给定速度—车体前进”完成仿真的基础。𝜔𝑓𝑙 , 𝜔𝑓𝑟 , 𝜔𝑟𝑙 , 𝜔𝑟𝑟
此外，仿真中还必须考虑关节阻尼和摩擦。前期调试时，模型生成后容易
出现轻微自溜、转向桥残余摆动或回中不干净等现象，因此需要在关节层加入
适当阻尼和摩擦参数[13-14,16]。这些参数虽然不写入运动学方程，却直接影







响车辆能否稳定停住、转向能否保持在目标角度，也使理论模型与仿真接口联
系起来。𝑐𝑑 𝑐𝑓
\subsection{本章小结}
本章针对机场飞机自动牵引场景，建立了双桥转向牵引车的平面运动学模